\subsection{Filtered groups}

Recall (Set Theory, III, p.~145) that an ordered set G is right (resp. left) filtered \(^{1}\) if for each pair \((x, y)\) of elements of G there exists \(z \in G\) such that \(x \leqslant z\) and \(y \leqslant z\) (resp. \(z \leqslant x\) and \(z \leqslant y\) ). Every right filtered ordered group G is also left filtered and conversely: indeed, since there exists \(z \in G\) such that \(-x \leqslant z\) and \(-y \leqslant z\) , we have \(-z \leqslant x\) and \(-z \leqslant y\) (VI, p.~3, Cor.). We will therefore speak simply of filtered groups.

PROPOSITION 4. --- For an ordered group G to be filtered it is necessary and sufficient that it be generated by its positive elements, that is that every element be the difference of two positive elements.

Indeed if G is filtered, then for each \(x \in G\) there exists a positive element z such that \(x \leqslant z\) , and x is the difference of the two positive elements z and z - x. Conversely, if x = u - v and y = w - t with u, v, w, t positive, then the element \(u + w\) is greater than x and greater than y.

PROPOSITION 5. --- If \((x_i)\) is a finite family of elements of a filtered group G, then there exists \(z \in G\) such that \(x_i + z\) is positive for each i.

If \(x_{i} = u_{i} - v_{i}\) , with \(u_{i}\) and \(v_{i}\) positive, it is enough to take \(z\) to be the sum of the family \((v_{i})\) .

